\documentclass{article}

\usepackage{amsmath,amssymb}
\usepackage{xurl}
\usepackage[hidelinks]{hyperref}

% Shared commands for notes in docs/paper-gaps/.

\DeclareUnicodeCharacter{2080}{\ifmmode{}_{0}\else\textsubscript{0}\fi}
\DeclareUnicodeCharacter{2081}{\ifmmode{}_{1}\else\textsubscript{1}\fi}
\DeclareUnicodeCharacter{2082}{\ifmmode{}_{2}\else\textsubscript{2}\fi}
\DeclareUnicodeCharacter{2099}{\ifmmode{}_{n}\else\textsubscript{n}\fi}

\newcommand{\Lean}{\textsc{Lean}}
\ExplSyntaxOn
\NewDocumentCommand{\leanid}{m}{
  \ifmmode
    \textsf{\detokenize{#1}}
  \else
    \group_begin:
    \urlstyle{sf}
    \tl_set:Nn \l_tmpa_tl {#1}
    \tl_replace_all:Nnn \l_tmpa_tl {\_} {_}
    \exp_args:NV \path \l_tmpa_tl
    \group_end:
  \fi
}
\ExplSyntaxOff
\providecommand{\tr}{\operatorname{tr}}
\newcommand{\tnleanIssueBase}{https://github.com/LionSR/TNLean/issues/}
\newcommand{\tnleanPrBase}{https://github.com/LionSR/TNLean/pull/}
\newcommand{\tnleanDocBase}{https://sirui-lu.com/TNLean/docs/}
\newcommand{\ghissue}[1]{\href{\tnleanIssueBase#1}{GitHub issue \##1}}
\newcommand{\ghpr}[1]{\href{\tnleanPrBase#1}{PR \##1}}
\newcommand{\doclink}[1]{\href{\tnleanDocBase#1.html}{\path{#1}}}

% Dirac notation and the identity operator, as in blueprint/src/macros/common.tex.
% \providecommand keeps note-local definitions authoritative.
\usepackage{dsfont}
\providecommand{\ket}[1]{|#1\rangle}
\providecommand{\bra}[1]{\langle #1|}
\providecommand{\braket}[2]{\langle #1 | #2 \rangle}
\providecommand{\ketbra}[2]{|#1\rangle\!\langle #2|}
\providecommand{\Id}{\mathds{1}}
\providecommand{\one}{\mathds{1}}
\providecommand{\C}{\mathbb{C}}
\providecommand{\MN}[1]{M_{#1}(\C)}
\providecommand{\MD}{M_D(\C)}

% External referencing for paper sources.  The build helper writes sanitized
% auxiliary files under build/paper-aux/ and excludes duplicated source labels.
\usepackage{xr}

% Import a generated paper auxiliary file with a caller-supplied label prefix.
% The candidate paths support builds from docs/paper-gaps/, the repository root,
% and build/paper-aux/ itself.
\newcommand{\importpaperaux}[2]{%
  \IfFileExists{../../build/paper-aux/#2.aux}{%
    \externaldocument[#1]{../../build/paper-aux/#2}%
  }{%
    \IfFileExists{build/paper-aux/#2.aux}{%
      \externaldocument[#1]{build/paper-aux/#2}%
    }{%
      \IfFileExists{#2.aux}{%
        \externaldocument[#1]{#2}%
      }{}%
    }%
  }%
}

% General external reference macros
\newcommand{\paperref}[2]{\ref{#1-#2}}
\newcommand{\papereqref}[2]{(\ref{#1-#2})}

% CPSV16 source (1606.00608 / MPDO-22-12-17-2.tex) external document and shortcuts
\importpaperaux{cpsv16-}{1606.00608/MPDO-22-12-17-2-tnlean}
\newcommand{\cpsvref}[1]{\paperref{cpsv16}{#1}}
\newcommand{\cpsveqref}[1]{\papereqref{cpsv16}{#1}}

% Verdict marker: \gapnote{<kind>}{<status>}. Kind is the policy
% classification (clarification, local-correction, scope-restriction,
% unfaithful, false-source, open-gap); status is open, wip (elimination
% underway), resolved, or historical. Machine-read by the site index;
% prints a verdict line.
\newcommand{\gapnote}[2]{\par\noindent\textbf{Verdict:} #1 (#2).\par}


\title{Historical Note: the Fixed-Block Cancellation Step in CPSV16 Theorem II.1}
\author{}
\date{2026-05-11}

\begin{document}
\maketitle

\section{Status}

\textbf{Historical proof archaeology, not an active gap.}
The corrected proportional fundamental theorem is
\leanid{MPSTensor.fundamentalTheorem_proportional_canonicalForm}.  It compares
active BNT families whose listed copy weights are nonzero and obtains a
bijection of the basis sectors with gauge-phase equivalences.  The unrestricted
statement printed in CPSV16 remains false: its weak BNT convention allows an
unrelated normal representative whose coefficient is identically zero, so the
claimed equality of BNT cardinalities fails.  The fixed-block cancellation
sentence therefore belongs only to the history of the superseded proof route.
The analysis below records why that route was abandoned.

\section{Source passage}

This note concerns the proof of CPSV16 Theorem~\texttt{thm1},
Appendix MPV proof line~1182 of \cite{Cirac2016MPDO_arXiv}. The paper fixes a block
$k_0\in\{1,\ldots,r_B\}$ and writes
\begin{align}
    \forall j, 
    \langle V^{(N)}(B_{k_0}),V^{(N)}(A_j)\rangle
    \xrightarrow[N\to\infty]{} 0
\end{align}
as a contradiction to eventual proportionality of the two full MPV families.
Under the BNT decomposition, the proportionality hypothesis is
projective proportionality: the scalar $c_N$ is nonzero at the lengths used in
the comparison, as recorded in
\path{docs/paper-gaps/cpsv16_nonzero_proportionality_reading.tex}. Thus
\begin{align}
    \sum_j \mu_j^N\,V^{(N)}(A_j)=c_N\sum_k \nu_k^N\,V^{(N)}(B_k),
    \qquad c_N\neq 0,
\end{align}
and isolating the selected block gives
\begin{align}
    \sum_j \mu_j^N\,V^{(N)}(A_j)-c_N\nu_{k_0}^N\,V^{(N)}(B_{k_0})
    =c_N\sum_{k\neq k_0}\nu_k^N\,V^{(N)}(B_k).
\end{align}
Corollary~\texttt{Lem1} in \cite{Cirac2016MPDO_arXiv} says that, from the left-hand
decay, the family
\begin{align}
    \dim\operatorname{span}
    \left(
      \left\{V^{(N)}(A_j)\right\}_{j=1}^{r_A}
      \cup
      \left\{V^{(N)}(B_{k_0})\right\}
    \right)
    =r_A+1
\end{align}
for all sufficiently large $N$. Equivalently, the displayed family is linearly
independent, so the selected block cannot be absorbed by the others without
contradiction. \cite[Appendix MPV proof, line~1182]{Cirac2016MPDO_arXiv}

\section{Formal boundary}

The formal boundary in \cite{Cirac2016MPDO_arXiv} is narrower:
\begin{align}
    \forall^\mathrm{cof} N, 
    \dim\operatorname{span}
    \left(
      \left\{V^{(N)}(A_j)\right\}_{j=1}^{r_A}
      \cup
      \left\{V^{(N)}(B_{k_0})\right\}
    \right)
    =r_A+1.
\end{align}
The surviving \leanid{SectorBNT} API does not contain this literal
single-block linear-independence formulation.  It supplies the stronger
full-family mechanism
\leanid{MPSTensor.IsBNTCanonicalForm.combined_family_eventually_li},
under all-pairwise cross-overlap decay.  That mechanism is sufficient for the
coefficient-comparison contradiction used by the corrected theorem, without
formalizing the source sentence as a separate fixed-block lemma.
The historical target was
\begin{align}
    (\forall j,  \langle V^{(N)}(B_{k_0}),V^{(N)}(A_j)\rangle\to 0)
    \Longrightarrow
    \neg\left(
      \sum_j \mu_j^N\,V^{(N)}(A_j)
      =
      c_N\sum_k \nu_k^N\,V^{(N)}(B_k)
    \right).
\end{align}
No current Lean declaration carries the retired fixed-right or fixed-left
one-copy fixed-block names.  Reintroducing this target would reconstruct an
abandoned proof architecture rather than close a current theorem.

The earlier conditional helpers requiring combined-family linear independence
(\path{_phase_sum_li}) and the residual-span variants have been retired in
the 2026-05-13 audit; see
\path{docs/audits/2026-05-13_cpsv16_ft_paper_vs_code_structural_map.md}.
Those helpers should not be used as replacements for the fixed-block sentence at
Appendix MPV proof line~1182, because the needed combined-family independence is not
supplied by Corollary~\texttt{Lem1}.  This mismatch explains why the fixed-block
route was retired; it is not a pending obligation on the corrected proportional
theorem.

Retirement note: this note intentionally no longer lists the retired no-tail,
leading-only, partner-identification, leading-tail, residual-span, or
\path{_phase_sum_li} declarations.  They were deleted as wrong-direction or
orphan scaffolding in the same audit.

\section{Retired discharge plan}

The plan recorded in
\path{docs/audits/2026-05-13_cpsv16_ft_sorry_discharge_plan.md} proposed proving
an exact eventual leading-coefficient identity
\begin{align}
    c_N\,\mu^B_{0}{}^N\,\zeta^N=(\mu^A_{0})^N
\end{align}
for all sufficiently large $N$ from the proportionality identity and the BNT
linear independence on the CPSV fixed-block family.  The proposal would then
have cloned the equal-MPV tail subtraction and induction architecture without
assuming linear independence of all residual blocks on both sides.  This plan
was superseded by direct coefficient comparison on active BNT families.

\subsection{Factored analytic discharge and its lower-bound obstruction}

\textit{Classification: historical analytic obstruction.}

The proposed factored analytic discharge used four explicit analytic
hypotheses:
\begin{enumerate}
  \item unit-modulus sector weights on both the $A$-family and the $B$-family;
  \item off-diagonal cross-overlap decay on the fixed-block family, i.e.,
        $\langle V^{(N)}(B_{k_0}), V^{(N)}(B_k)\rangle \to 0$ for $k \neq k_0$;
  \item a nonzero limit of the self-overlap at the fixed block,
        $\langle V^{(N)}(B_{k_0}), V^{(N)}(B_{k_0})\rangle \to \ell \neq 0$; and
  \item an eventual positive lower bound $\|c_N\| \geq \delta > 0$.
\end{enumerate}
Hypotheses~(1)--(3) are present in the source via the BNT structure
\cite[Theorem~\texttt{thm1}]{Cirac2016MPDO_arXiv}.  Hypothesis~(4) is not
directly supplied: the cited argument rules out cancellation by appeal to the
dominant-adjusted scalar limit, but the retired plan did not derive the explicit
lower bound $\|c_N\|\geq\delta$ from the canonical-form data.  The earlier
non-cancellation helpers carrying this lower-bound hypothesis belonged to the
retired Tier~3 route and should not be cited as present declarations.  No
current theorem depends on this lower-bound argument.

\bibliographystyle{alpha}
\bibliography{references}

\end{document}
